\documentclass{article}
%DIF LATEXDIFF DIFFERENCE FILE
%DIF DEL build/arxiv_changes/main_ml4ps.tex   Sun Sep 27 14:38:24 2026
%DIF ADD main.tex                             Sun Sep 27 14:03:18 2026

%DIF 3c3
%DIF < \usepackage{neurips_2026}
%DIF -------
\usepackage[preprint]{neurips_2026}
 %DIF > 
%DIF -------

\usepackage[utf8]{inputenc}
\usepackage[T1]{fontenc}
\usepackage{hyperref}
\usepackage{url}
\usepackage{booktabs}
\usepackage{amsfonts}
\usepackage{amsmath}
\usepackage{amssymb}   % \lesssim in Appendix A; amsmath alone does not define it
\usepackage{nicefrac}
\usepackage{microtype}
\usepackage{xcolor}
\usepackage{graphicx}
\usepackage{float}     % [H] for the Appendix C figure, so it stays with its paragraph
%DIF 18-20d18
%DIF < \setlength{\abovecaptionskip}{4pt}   % float spacing tightened to hold the body on four pages
%DIF < \setlength{\textfloatsep}{10pt}
%DIF < \setlength{\floatsep}{8pt}
%DIF -------

%DIF 22-29d19
%DIF < % The workshop prescribes the first-page footer text. The style file is otherwise untouched.
%DIF < \makeatletter
%DIF < \renewcommand{\@noticestring}{%
%DIF <   Submitted to the 9th Workshop on Machine Learning and the Physical Sciences (ML4PS 2026).
%DIF <   Do not distribute.%
%DIF < }
%DIF < \makeatother
%DIF < 
%DIF -------
\title{What does self-supervision add to engineered features\\for stellar light curves?}

%DIF 32a21-24
% Author block (roadmap 2.1). Order matches the OpenReview record (Zijun Liu, Yue Ma, Christopher
 %DIF > 
% Theissen); Yue Ma's affiliation is her OpenReview profile's current entry. Co-author e-mail
 %DIF > 
% addresses are not printed without their consent. "Hal\i c\i o\u{g}lu" is spelled with the escapes
 %DIF > 
% rather than raw UTF-8 so the file stays ASCII.
 %DIF > 
%DIF -------
\author{\DIFdelbegin \DIFdel{Anonymous Author(s)}\DIFdelend %DIF > 
  \DIFaddbegin \DIFadd{Zijun Liu}\DIFaddend \\
  \DIFdelbegin \DIFdel{Affiliation}\DIFdelend \DIFaddbegin \DIFadd{University of California, San Diego}\DIFaddend \\
  \DIFdelbegin \DIFdel{Address}\DIFdelend \DIFaddbegin \DIFadd{\texttt{brl055@ucsd.edu}}\DIFaddend \\
  \DIFdelbegin \DIFdel{\texttt{email}}\DIFdelend \DIFaddbegin \DIFadd{\texttt{liuzijun6688@gmail.com}
  }\And
  \DIFadd{Yue Ma}\\
  \DIFadd{Hal\i c\i o\u{g}lu Data Science Institute}\\
  \DIFadd{University of California, San Diego
  }\And
  \DIFadd{Christopher A. Theissen}\\
  \DIFadd{Department of Astronomy and Astrophysics}\\
  \DIFadd{University of California, San Diego
}\DIFaddend }

%DIF PREAMBLE EXTENSION ADDED BY LATEXDIFF
%DIF UNDERLINE PREAMBLE %DIF PREAMBLE
\RequirePackage[normalem]{ulem} %DIF PREAMBLE
\RequirePackage{color}\definecolor{RED}{rgb}{1,0,0}\definecolor{BLUE}{rgb}{0,0,1} %DIF PREAMBLE
\providecommand{\DIFaddtex}[1]{{\protect\color{blue}\uwave{#1}}} %DIF PREAMBLE
\providecommand{\DIFdeltex}[1]{{\protect\color{red}\sout{#1}}} %DIF PREAMBLE
%DIF SAFE PREAMBLE %DIF PREAMBLE
\providecommand{\DIFaddbegin}{} %DIF PREAMBLE
\providecommand{\DIFaddend}{} %DIF PREAMBLE
\providecommand{\DIFdelbegin}{} %DIF PREAMBLE
\providecommand{\DIFdelend}{} %DIF PREAMBLE
\providecommand{\DIFmodbegin}{} %DIF PREAMBLE
\providecommand{\DIFmodend}{} %DIF PREAMBLE
%DIF FLOATSAFE PREAMBLE %DIF PREAMBLE
\providecommand{\DIFaddFL}[1]{\DIFadd{#1}} %DIF PREAMBLE
\providecommand{\DIFdelFL}[1]{\DIFdel{#1}} %DIF PREAMBLE
\providecommand{\DIFaddbeginFL}{} %DIF PREAMBLE
\providecommand{\DIFaddendFL}{} %DIF PREAMBLE
\providecommand{\DIFdelbeginFL}{} %DIF PREAMBLE
\providecommand{\DIFdelendFL}{} %DIF PREAMBLE
%DIF HYPERREF PREAMBLE %DIF PREAMBLE
\providecommand{\DIFadd}[1]{\texorpdfstring{\DIFaddtex{#1}}{#1}} %DIF PREAMBLE
\providecommand{\DIFdel}[1]{\texorpdfstring{\DIFdeltex{#1}}{}} %DIF PREAMBLE
\providecommand{\DIFscaledelfig}{0.5}
%DIF HIGHLIGHTGRAPHICS PREAMBLE %DIF PREAMBLE
\RequirePackage{settobox} %DIF PREAMBLE
\RequirePackage{letltxmacro} %DIF PREAMBLE
\newsavebox{\DIFdelgraphicsbox} %DIF PREAMBLE
\newlength{\DIFdelgraphicswidth} %DIF PREAMBLE
\newlength{\DIFdelgraphicsheight} %DIF PREAMBLE
% store original definition of \includegraphics %DIF PREAMBLE
\LetLtxMacro{\DIFOincludegraphics}{\includegraphics} %DIF PREAMBLE
\providecommand{\DIFaddincludegraphics}[2][]{{\color{blue}\fbox{\DIFOincludegraphics[#1]{#2}}}} %DIF PREAMBLE
\providecommand{\DIFdelincludegraphics}[2][]{% %DIF PREAMBLE
\sbox{\DIFdelgraphicsbox}{\DIFOincludegraphics[#1]{#2}}% %DIF PREAMBLE
\settoboxwidth{\DIFdelgraphicswidth}{\DIFdelgraphicsbox} %DIF PREAMBLE
\settoboxtotalheight{\DIFdelgraphicsheight}{\DIFdelgraphicsbox} %DIF PREAMBLE
\scalebox{\DIFscaledelfig}{% %DIF PREAMBLE
\parbox[b]{\DIFdelgraphicswidth}{\usebox{\DIFdelgraphicsbox}\\[-\baselineskip] \rule{\DIFdelgraphicswidth}{0em}}\llap{\resizebox{\DIFdelgraphicswidth}{\DIFdelgraphicsheight}{% %DIF PREAMBLE
\setlength{\unitlength}{\DIFdelgraphicswidth}% %DIF PREAMBLE
\begin{picture}(1,1)% %DIF PREAMBLE
\thicklines\linethickness{2pt} %DIF PREAMBLE
{\color[rgb]{1,0,0}\put(0,0){\framebox(1,1){}}}% %DIF PREAMBLE
{\color[rgb]{1,0,0}\put(0,0){\line( 1,1){1}}}% %DIF PREAMBLE
{\color[rgb]{1,0,0}\put(0,1){\line(1,-1){1}}}% %DIF PREAMBLE
\end{picture}% %DIF PREAMBLE
}\hspace*{3pt}}} %DIF PREAMBLE
} %DIF PREAMBLE
\LetLtxMacro{\DIFOaddbegin}{\DIFaddbegin} %DIF PREAMBLE
\LetLtxMacro{\DIFOaddend}{\DIFaddend} %DIF PREAMBLE
\LetLtxMacro{\DIFOdelbegin}{\DIFdelbegin} %DIF PREAMBLE
\LetLtxMacro{\DIFOdelend}{\DIFdelend} %DIF PREAMBLE
\DeclareRobustCommand{\DIFaddbegin}{\DIFOaddbegin \let\includegraphics\DIFaddincludegraphics} %DIF PREAMBLE
\DeclareRobustCommand{\DIFaddend}{\DIFOaddend \let\includegraphics\DIFOincludegraphics} %DIF PREAMBLE
\DeclareRobustCommand{\DIFdelbegin}{\DIFOdelbegin \let\includegraphics\DIFdelincludegraphics} %DIF PREAMBLE
\DeclareRobustCommand{\DIFdelend}{\DIFOaddend \let\includegraphics\DIFOincludegraphics} %DIF PREAMBLE
\LetLtxMacro{\DIFOaddbeginFL}{\DIFaddbeginFL} %DIF PREAMBLE
\LetLtxMacro{\DIFOaddendFL}{\DIFaddendFL} %DIF PREAMBLE
\LetLtxMacro{\DIFOdelbeginFL}{\DIFdelbeginFL} %DIF PREAMBLE
\LetLtxMacro{\DIFOdelendFL}{\DIFdelendFL} %DIF PREAMBLE
\DeclareRobustCommand{\DIFaddbeginFL}{\DIFOaddbeginFL \let\includegraphics\DIFaddincludegraphics} %DIF PREAMBLE
\DeclareRobustCommand{\DIFaddendFL}{\DIFOaddendFL \let\includegraphics\DIFOincludegraphics} %DIF PREAMBLE
\DeclareRobustCommand{\DIFdelbeginFL}{\DIFOdelbeginFL \let\includegraphics\DIFdelincludegraphics} %DIF PREAMBLE
\DeclareRobustCommand{\DIFdelendFL}{\DIFOaddendFL \let\includegraphics\DIFOincludegraphics} %DIF PREAMBLE
%DIF AMSMATHULEM PREAMBLE %DIF PREAMBLE
\makeatletter %DIF PREAMBLE
\let\sout@orig\sout %DIF PREAMBLE
\renewcommand{\sout}[1]{\ifmmode\text{\sout@orig{\ensuremath{#1}}}\else\sout@orig{#1}\fi} %DIF PREAMBLE
\makeatother %DIF PREAMBLE
%DIF COLORLISTINGS PREAMBLE %DIF PREAMBLE
\RequirePackage{listings} %DIF PREAMBLE
\RequirePackage{color} %DIF PREAMBLE
\lstdefinelanguage{DIFcode}{ %DIF PREAMBLE
%DIF DIFCODE_UNDERLINE %DIF PREAMBLE
  moredelim=[il][\color{red}\sout]{\%DIF\ <\ }, %DIF PREAMBLE
  moredelim=[il][\color{blue}\uwave]{\%DIF\ >\ } %DIF PREAMBLE
} %DIF PREAMBLE
\lstdefinestyle{DIFverbatimstyle}{ %DIF PREAMBLE
	language=DIFcode, %DIF PREAMBLE
	basicstyle=\ttfamily, %DIF PREAMBLE
	columns=fullflexible, %DIF PREAMBLE
	keepspaces=true %DIF PREAMBLE
} %DIF PREAMBLE
\lstnewenvironment{DIFverbatim}[1][]{\lstset{style=DIFverbatimstyle}}{} %DIF PREAMBLE
\lstnewenvironment{DIFverbatim*}[1][]{\lstset{style=DIFverbatimstyle,showspaces=true}}{} %DIF PREAMBLE
\lstset{extendedchars=\true,inputencoding=utf8}

%DIF END PREAMBLE EXTENSION ADDED BY LATEXDIFF

\begin{document}

\maketitle

\begin{abstract}
TESS light curves encode a star's pulsation, rotation, binarity and oscillations, among other physical
properties, but the survey
produces far more light curves than can be labelled by hand. Astronomers have engineered a few
dozen physically motivated features from them, which together remain a strong baseline. Previous work has trained
encoders on light curves and judged them by their own downstream scores, but has not asked what
they add to the engineered features. We pre-train a convolutional encoder on unlabelled TESS light curves to predict each
window's latent representation from those of its neighbours. We then freeze it, add its features to
the 25 engineered features, and read the combination out with one fixed linear probe on ten tasks:
five variability classes, four asteroseismic targets and rotation period. Adding the learned features
improves the readout on 7 of 10 tasks over the engineered features alone, while an untrained encoder
does not. The combination also outperforms two networks trained end-to-end on the same labelled
stars on most tasks. An encoder trained without labels therefore gives astronomers new features
that complement the ones they already use rather than replace them.
\end{abstract}

\section{Introduction}

A star's light curve carries its pulsation, rotation, eclipses and oscillations, and with them its
age, interior and companions. Wide-field surveys now produce light curves far faster
than they can be labelled. TESS
\citep{ricker2015} returns 2-minute or 20-second \citep{huber2022} light curves for hundreds of
thousands of pre-selected stars, and full-frame images of every star in the field at cadences that
change between mission cycles (30 minutes, then 10, now 200 seconds). LSST \citep{ivezic2019}
will add \emph{billions} more light curves, albeit at a cadence of days. Self-supervised learning is the
natural response: pre-train one encoder on the unlabelled archive, then attach a cheap readout for
each task. Encoders of this shape exist for ground-based, Gaia and Kepler light curves
\citep{donoso2025astromer2,zuo2025falco,rui2026jepa} and are evaluated by classification accuracy,
sometimes against engineered features \citep{rui2026jepa}; a label-free TESS representation
\citep{poznanski2026} is evaluated by whether repeat observations of one star lie close in its embedding space.
None measures what they add to engineered features.

In this domain, the baseline that a learned encoder must improve on is already strong. Decades of variable-star
research have distilled light curves into compact engineered summaries: periodogram peak structure,
band-limited power, spectral entropy and robust scatter statistics. Pipelines built on them remain
highly competitive. The benchmark of \citet{li2025staremb} makes the point
sharply: across seven expert-labelled variable-star classes, no time-series foundation model they
tested strictly surpassed the long-established domain feature baseline on classification. A learned
encoder for this data therefore has to answer one question: do its features add predictive signal to
the engineered features under the same readout?

We measure that directly. We pre-train an encoder on TESS light curves, freeze it, and add its
features to the 25 engineered features under one fixed linear readout. We ask two things: what the
learned features add, and whether the combination, which trains only a linear model on the labels,
outperforms two networks trained end-to-end on the same labels. Our contributions are:

\begin{itemize}
\item \textbf{The learned features add to the engineered features rather than replacing them.}
Adding the learned features to the engineered ones improves a fixed linear readout beyond
$2\,\mathrm{SE}$ on 7 of 10 tasks; an untrained encoder gives no such gain, so the gain comes from pre-training,
not from merely adding 128 columns.
\item \textbf{The combined pipeline, read out by a fixed linear model, outperforms the two
supervised networks we train end-to-end on the same labels on most tasks.} Against a supervised Conv1D with the
encoder's architecture it is ahead on 6
of 10 tasks and indistinguishable on the other 4; against an MLP on raw flux it is ahead on 9 of 10.
\end{itemize}

\section{Data and method}

\paragraph{Data.} We use 2-minute SPOC PDCSAP light curves of stars brighter than TESS magnitude
$T = 10$ from Sectors 27--101. Pre-training and five of the ten probes (pulsating, eclipsing binary,
rotation, transit and rotation period) use a label-enriched subset of 13{,}470 stars, split 70/15/15
by star: every star labelled as a transit host, eclipsing binary or pulsator that has a usable
window, plus a random sample of 10{,}000 quiet stars, listed in none of our variability catalogues. The pre-training
loss never sees a label, but this pool is label-enriched, not a random draw from the survey. The other five probes draw their labels from catalogues that barely overlap this subset and are scored on a second pool of 22{,}860 stars built the same way (every catalogue positive
with a usable window plus 10{,}000 quiet stars, split 70/15/15). No test star of the first pool
enters pre-training. \DIFaddbegin \DIFadd{The first pool admits positives only through the transit, eclipsing-binary and pulsating catalogues and screens its quiet stars against the rotation catalogue as well, so every rotator in it carries a second label and it holds no rotation-only star; Appendix~J rescores the two rotation tasks on a third pool of 106}{\DIFadd{,}}\DIFadd{284 stars, the corpus outside the first pool, where rotators are a natural 13\%. }\DIFaddend About 8\% of the second pool's test stars (265 of 3{,}429) were in the pre-training subset; removing them leaves all seven gains in place
(Appendix~I). Each light curve is split into contiguous segments wherever the time axis
breaks, at a missing cadence or a jump longer than five median cadences; segments are
median-normalised by their median absolute deviation and sliced into non-overlapping windows of 256
cadences ($256 \times 120$\,s $\approx 8.5$\,h); restricting to one cadence keeps every window the
same span of time. Windows containing a missing value are discarded, never interpolated or padded,
and segments are never stitched across sector or gap boundaries. Pre-training consumes
sequences of exactly 16 consecutive windows ($\approx 5.69$\,d), cut at a random offset within a
segment, and uses no labels.

\paragraph{Encoder.} A 1D convolutional variational encoder \citep{kingma2014} maps each window to a
128-dimensional Gaussian posterior, and a GRU \citep{cho2014} predicts the next window's latent from
its predecessors, run both forward and backward over the sequence. The objective is
\begin{align*}
\mathcal{L} = {} & \|\hat{x}-x\|^{2} + \beta\,\mathrm{KL}
+ \lambda\big[\|\hat{z}_{t+1}-\bar{z}_{t+1}\|^{2} + \|\hat{z}_{t-1}-\bar{z}_{t-1}\|^{2}\big] \\
& + w\big[\|\log P_{H}(\hat{x})-\log P_{H}(x)\|^{2} + \|\Delta\hat{x}-\Delta x\|^{2}\big],
\end{align*}
where every term is a mean squared error, $\bar{z}$ the stop-gradient target of the forward and
backward predictors, $P_{H}$ the Hann-tapered power spectrum, $\Delta$ the first difference, and
$(\beta,\lambda,w) = (0.1, 60, 0.3)$, $\beta$ warmed up linearly over ten epochs. The encoder is trained with six seeds and then frozen; no gradient from any label reaches
it.

\paragraph{Representations and readout.} A star's learned features are $\mu$, the mean, over every
window of its first segment (16 for 95\% of stars, 20 for the rest), of the per-window posterior
means. The engineered baseline is a fixed 25-feature vector modelled on the T'DA variability pipeline
\citep{audenaert2021}, computed on the same first segment, 19 from its periodogram and 6 time-domain statistics; Appendix~H defines each feature. Readouts are fixed in advance and never
tuned on a task, with inputs standardised using training-split statistics alone: class-balanced $L_{2}$ logistic regression
at $C = 1$ for the detection tasks, ridge for the two regressions with its penalty set by leave-one-out
cross-validation. We compare four
arms, each the same readout on a different feature set: \texttt{features}, $\mu$,
\texttt{features}\,$\oplus\,\mu$, and, as a control, the same fusion built on an \emph{untrained}
encoder; the two supervised baselines below are the fifth and sixth arms.
Appendix~D repeats the first three under XGBoost.

\paragraph{Supervised baselines.} Two end-to-end baselines see exactly the same input windows, star
sets and splits. \textbf{C1} reuses the encoder's convolutional trunk unchanged, adds a 128-unit linear layer
and a scalar output, and is trained per task with a class-weighted loss on a star-level score, the mean of the star's window
logits computed inside the network so that the gradient flows through the mean. \textbf{C2} replaces the convolutional trunk with a
dense one and changes nothing else, so C1 versus C2 isolates the convolutional inductive bias. Both
use one pre-registered regularisation setting for all tasks, select on a validation split no probe
has ever touched, and are trained with six seeds.

\paragraph{Tasks and metrics.} We score ten tasks with one probe per physical quantity: pulsating
variables \citep{gao2025}, eclipsing binaries \citep{prsa2022}, rotation and rotation period
\citep{boyle2026}, transits \citep{guerrero2021}, flare activity \citep{seli2025,vida2021},
oscillating red giants and their $\nu_{\max}$ \citep{hon2021}, solar-like oscillators
\citep{hatt2023}, and RGB versus helium-burning discrimination \citep{sreenivas2026}. Here
$\nu_{\max}$ is the frequency of maximum oscillation power. Detection tasks are scored by the area
under the precision--recall curve (PR-AUC), whose floor is the positive-class prevalence rather than
$0.5$; prevalence is therefore quoted with every detection score. The two regressions are scored by
$R^{2}$.

\paragraph{Uncertainties.} All error bars are $2\times$ the standard error of the mean over seeds.
For the $\mu$ arms the seeds are independent pre-training runs (six); for C1 and C2 they are
initialisation and shuffling seeds (six); the engineered arm is
deterministic. The seeds of the $\mu$ arms and of the supervised arms are unrelated, so differences
between them are treated as unpaired and their bands combined as $\sqrt{\mathrm{SE}_a^{2}+\mathrm{SE}_b^{2}}$.
A difference counts only when it exceeds the combined $2\,\mathrm{SE}$ band. These bands
measure re-training spread only; a paired bootstrap over the test stars (1{,}000 resamples, seeds
averaged; Appendix~G) checks that no verdict hinges on the particular test stars.

\section{Results}

\begin{table}[t]
\caption{Ten-task scorecard. Frozen arms: a fixed logistic or ridge readout, or XGBoost on the same
inputs (Appendix~D). Bold: best point estimate of the three arms under each frozen
readout\DIFaddbeginFL \DIFaddFL{, with no significance test attached; bold in the appendix tables instead marks a difference
resolved beyond its error bar, so the two conventions are not interchangeable}\DIFaddendFL . $n$: test stars;
prev.: positive-class prevalence.}
\label{tab:scorecard}
\centering
\scriptsize
\input{tables/table1_scorecard}
\end{table}

\begin{figure}[t]
\centering
\includegraphics[width=\linewidth]{figures/fig1_deltas_xgb.pdf}
\caption{The two claims. \textbf{(a)} What the learned features add to the engineered features, under
the fixed linear readout of Table~\ref{tab:scorecard} (solid) and under an XGBoost readout on the
same inputs (faded; Appendix~D). \textbf{(b)} The fusion probe against the two
supervised baselines. Bars are $2\,\mathrm{SE}$.}
\label{fig:deltas}
\end{figure}

\paragraph{The learned features add to the engineered features.} Table~\ref{tab:scorecard} and
Figure~\ref{fig:deltas}a give the comparison. Adding $\mu$ to the 25 features improves the
fixed readout beyond $2\,\mathrm{SE}$ on 7 of 10 tasks, led by solar-like oscillators $+0.071$,
pulsating $+0.057$ and flare activity $+0.051$, and reaching $+0.013$ ($R^{2}$) on rotation period.
Rotation is positive but just short of the
bar ($+0.016$ against $0.016$). \DIFaddbegin \DIFadd{On the rotation pool of Appendix~J, where the two classes differ in rotation and nothing else, the rotation gain is $+0.055 \pm 0.006$ but an untrained encoder also gains $+0.029$, so about half of it is what any 128 extra columns buy on 74}{\DIFadd{,}}\DIFadd{000 training stars; the rotation-period gain there is $+0.034 \pm 0.002$ with the untrained control at zero. Neither verdict moves the count. }\DIFaddend The two negative tasks are the ones where 128 extra columns cost more
than they add: the 25 features already cover the task (oscillating giants, $-0.003$ on a baseline
of $0.920$) or the training set is too small to fit 128 more weights (RGB versus helium-burning, $-0.050$ on
755 training stars). On both, an untrained encoder loses by about the same amount as the trained
fusion, within the error band; that is what redundant columns do to a linear fit, not a sign that
$\mu$ carries misleading content (Appendix~C). Across all ten tasks the untrained control is
negative on seven and never exceeds $+0.016$ (Table~D.1), so the gains are not what any 128 extra columns would
buy on this feature set; they are specific to what pre-training learned. $\mu$ on its own does not
replace the features (ahead on only 4 of 10 tasks alone), as \citet{li2025staremb} also find for
time-series foundation models; we therefore report fusion rather than substitution \citep{makinen2024}.

\paragraph{The fusion probe outperforms the two supervised networks on most tasks.} Against C1, which shares
the convolutional trunk, labels and input windows but is trained end-to-end, the fusion
probe is ahead on 6 of 10 tasks and indistinguishable from it on the other 4
(Figure~\ref{fig:deltas}b). Against
C2, an MLP on raw flux, it is ahead on 9 of 10 and behind on one (RGB versus helium-burning, 161
stars). C1 is not weak: it is ahead of the engineered arm on 5 of 10.
A paired star-bootstrap (Appendix~G) excludes zero on 5 of the 7 fusion gains, 5
of the 6 leads over C1 and 7 of the 9 over C2.
This holds for this architecture on 669 to 16{,}002 labelled stars per task, not for
supervision in general.

\paragraph{The GRU prediction term raises the gain on every task.} Removing the
GRU prediction term, all else fixed and the same six
seeds, reduces the fusion gain on all 10 tasks, and by more than $2\,\mathrm{SE}$ on 6 when the two
arms are differenced seed by seed (Table~D.1). Without it the gain falls by more
than half on 5 of the 8 tasks where fusion gains.

\section{Discussion}

\paragraph{Pick the encoder by a held-out probe, not by its loss.} Across the ten pre-training recipes of
one sweep, a lower validation reconstruction loss went with a \emph{lower} probe score: the Spearman
$\rho$ between negated loss and PR-AUC runs from $-0.58$ to $-0.84$ over four tasks
\DIFaddbegin \DIFadd{(Figure~\ref{fig:valloss_a})}\DIFaddend . The recipe with the
lowest loss is the one without the prediction term on all four, and choosing it would cost
$0.021$--$0.049$ PR-AUC\DIFaddbegin \DIFadd{. On the family our encoder belongs to the anti-correlation does not survive as
such --- it holds on two of the four tasks and reverses on two --- but the loss optimum is still never
the probe optimum }\DIFaddend (Appendix~E). \DIFaddbegin \DIFadd{The same objective carries a second cost: it develops
a localised artifact inside each window, and none of the interventions we tried on four axes cleared
both pre-registered acceptance criteria, so the recipe we ship does so by the stated fallback rule
rather than by winning (Appendix~F). }\DIFaddend An encoder here therefore has to be
chosen by a held-out probe, not by its own loss.

\DIFaddbegin \begin{figure}[t]
\centering
\includegraphics[width=0.85\linewidth]{figures/figB_valloss_a.pdf}
\caption{\DIFaddFL{Validation reconstruction loss against downstream probe score, one point per pre-training
recipe and one panel per task. The ten recipes share corpus, architecture and loss terms and differ
only in the GRU prediction term (off, forward, forward+backward, multi-step, the last three at three
weights), four seeds each. $x$: best post-warm-up validation reconstruction loss, averaged over the
recipe's seeds; $y$: PR-AUC of the fixed linear probe on $\mu$; $\rho$: Spearman's rank correlation
between negated loss and score. The loss-optimal recipe is the dynamics-free one (red) on all four
tasks, and it is the worst or next-to-worst probe on each. Appendix~E repeats the
measurement on the family this paper's encoder belongs to.}}
\label{fig:valloss_a}
\end{figure}

\DIFaddend \paragraph{\textbf{Limitations and next steps.}} \emph{Readout.} The four linear-readout arms (\texttt{features},
$\mu$, \texttt{features}\,$\oplus\,\mu$, and the untrained control) share one linear model. It has few
parameters, so it cannot rescue weak features, and it is identical across the four, so the comparison
is between feature sets. Our claim is confined to these 25 features and a linear readout: under
XGBoost on the same inputs the gain exceeds $2\,\mathrm{SE}$ on 4 of 10 tasks and is negative on one\DIFaddbegin \DIFadd{;
where the nonlinear readout raises the features-only score, the $\mu$ gain falls, with a rank
correlation of $0.92$ across the ten tasks }\DIFaddend (Appendix~D). \emph{Data.} Every arm,
the supervised baselines included, sees one 5.7-day segment of a bright 2-minute star. The fainter,
far more numerous stars in the full-frame images, and light curves spanning several sectors, are the
obvious next experiments. \emph{Labels.} Our
negatives are stars absent from the label catalogues, not stars verified to lack the signal, so
catalogue incompleteness adds positive-unlabelled label noise and the detection scores may be
conservative. The flare label marks stars that have flared somewhere in their TESS data, not a flare inside
the 5.7-day segment we encode, so the task is recognising a flaring star, not finding a flare. \emph{Label scarcity.} When
the readout is trained on a few percent of the labels, the advantage narrows on 7 of 10 tasks, so how
much pre-training helps when labels are scarce remains open. \emph{Discovery.} Every task here has a
catalogue; whether $\mu$ also separates kinds of variability that no catalogue labels yet, for instance
as clusters in the latent space, is untested. \DIFaddbegin \DIFadd{We leave it as an exploratory direction rather than a
result: with no labels there is nothing to score such a separation against, so it needs a design of
its own.

}\DIFaddend 

\paragraph{Reproducibility and disclosure.} \S2 and Appendix~H give the data,
objective, features and readouts; the encoder has four convolutional stages (widths $32/64/128/256$,
kernel 5, each halving length), a 128-d latent and a one-layer GRU of width 256. Code\DIFdelbegin \DIFdel{and score
tables will be released on publication}\DIFdelend \DIFaddbegin \DIFadd{, configuration
files and the score tables behind Table~\ref{tab:scorecard} and the appendix tables are available at
}\url{https://github.com/BrightonLiu-zZ/stellar-ssl/tree/arxiv-v1}\DIFaddend ; generative AI assisted with
analysis code and prose editing.


\clearpage
\bibliographystyle{plainnat}
\bibliography{refs}
\end{document}
